\chapter*{Vorwort}
\addcontentsline{toc}{chapter}{Vorwort}

\section*{Worum es geht}

Auf jeder Schokoladenpackung, in jedem Funkpaket, auf jeder
zerkratzten CD und in jedem Datamatrix-Code auf einem Versandetikett
arbeitet im Hintergrund \emph{Mathematik}, die Fehler erkennt und
korrigiert. Dieses Buch baut diese Mathematik Stück für Stück auf —
von der einfachsten Prüfziffer bis zum vollständigen Datamatrix-Code
mit Reed-Solomon.

\section*{Über Greta}

Das Buch ist als zweiwöchiges Praktikum gedacht und richtet sich
direkt an „Greta", eine Schülerin der zehnten Klasse. Greta gibt es
wirklich; aber selbst wenn du keine Schülerin namens Greta bist, kannst
du das Buch als Greta lesen. Der direkte Ton ist Absicht: niemand
erklärt etwas Schwieriges in der dritten Person.

Vorwissen, das wir voraussetzen: Schulmathematik bis Klasse 10
(Bruchrechnen, Funktionen, ein bisschen Beweise mitdenken) und etwas
Python-Erfahrung (\texttt{for}-Schleifen, Listen, Funktionen — oder die
Bereitschaft, das nebenbei aufzuholen).

\section*{Wie das Buch aufgebaut ist}

Jedes Kapitel ist ein \emph{Tag}. Ein Tag dauert ungefähr zwei Stunden
und besteht aus mehreren Blöcken mit Zeitangaben. In jedem Block
wechseln sich drei Aufgabentypen ab, die du am Strich und Symbol auf
der Seitenaußenkante erkennst:

\begin{itemize}
\item \textcolor{cBleistift}{\textbf{Bleistift-Übungen}} (orange) — mit
  Stift und Karopapier. Selbstdenken vor Selbstrechnen vor Selbstprüfen.
\item \textcolor{cPython}{\textbf{Python-Einheiten}} (blau) — Code in
  Thonny abtippen, ausführen, durchspielen. Computer brute-forcen
  schneller als wir; das ist eines unserer Lieblingsspielzeuge.
\item \textcolor{cWerkzeug}{\textbf{Werkzeug-Checks}} (grün) — kurze
  Setup-Schritte, damit du nicht beim Installieren von Thonny stecken
  bleibst.
\end{itemize}

Im Anhang stehen die \textbf{Lösungen} (Anhang A) und ein
\textbf{Glossar} aller Fachbegriffe. Lösungen sind absichtlich
ausführlich; den Aha-Moment hat man trotzdem nur, wenn man sie erst
nach eigenem Versuch liest.

\section*{Material}

\begin{itemize}
\item Bleistift, Radiergummi, kariertes Papier
\item drei Buntstifte für die Diagramme an Tag 3
\item ein Laptop mit Python (Thonny ist die empfohlene Umgebung)
\item ein Produkt mit Strichcode (Schokolade, Buch — egal) für Tag 1
\end{itemize}

\section*{Tempo}

Die Tagesabschnitte sind als Häppchen gedacht. Wenn ein Tag sich nach
zwei Stunden noch nicht zu Ende anfühlt, ist es kein Verbrechen, ihn
auf zwei Sitzungen aufzuteilen. Tag 5 (endliche Körper) ist
konzeptionell der schwerste — gerade dort darf Pause sein.

Viel Spaß beim Lesen, Rechnen und Tippen.
